\section{Ratios de rentabilidad}
Los ratios de rentabilidad evalúan la capacidad de generar beneficios respecto a ventas, activos o patrimonio. Son útiles para medir desempeño, pero deben compararse con periodos anteriores, sector y costo de capital.

\subsection{Margen neto}
\[
\text{Margen neto}=\frac{\text{Utilidad neta}}{\text{Ventas netas}}\times 100
\]
Datos: ventas netas \soles{1,500,000.00} y utilidad neta \soles{135,000.00}.
\[
\frac{135,000.00}{1,500,000.00}\times 100=9.00\text{ \char37}
\]
\textbf{Interpretación:} por cada \soles{100.00} de ventas, la empresa obtiene \soles{9.00} de utilidad neta.

\newpage
\subsection{ROA}
\[
\text{ROA}=\frac{\text{Utilidad neta}}{\text{Activo promedio}}\times 100
\]
Activo promedio: \soles{1,500,000.00}.
\[
\frac{135,000.00}{1,500,000.00}\times 100=9.00\text{ \char37}
\]
\textbf{Interpretación:} los activos generaron una rentabilidad neta de 9.00 \char37{}.

\newpage
\subsection{ROE}
\[
\text{ROE}=\frac{\text{Utilidad neta}}{\text{Patrimonio promedio}}\times 100
\]
Patrimonio promedio: \soles{900,000.00}.
\[
\frac{135,000.00}{900,000.00}\times 100=15.00\text{ \char37}
\]
\textbf{Interpretación:} el patrimonio generó una rentabilidad de 15.00 \char37{}. Un ROE mayor que el ROA puede relacionarse con apalancamiento, pero no siempre es favorable si la deuda es costosa.
